\documentclass[11pt]{article}

\usepackage[a4paper,margin=1in]{geometry}
\usepackage{amsmath,amssymb,amsthm}
\usepackage{mathtools}
\usepackage{microtype}
\usepackage{booktabs}
\usepackage{array}
\usepackage{graphicx}
\usepackage{tikz}
\usepackage{pgfplots}
\pgfplotsset{compat=1.18}
\usetikzlibrary{arrows.meta,positioning,calc,shapes.geometric,shapes.misc,fit,backgrounds,decorations.pathreplacing,patterns}
\usepackage[hidelinks]{hyperref}
\usepackage{xcolor}
\usepackage{enumitem}
\usepackage{caption}
\usepackage{tabularx}
\usepackage{verbatim}

\definecolor{accent}{HTML}{1f3a93}
\definecolor{accent2}{HTML}{c0392b}
\definecolor{accent3}{HTML}{27ae60}
\definecolor{rule}{HTML}{888888}
\definecolor{lightgray}{HTML}{dcdcdc}

\hypersetup{
  colorlinks=true,
  linkcolor=accent,
  citecolor=accent,
  urlcolor=accent,
}

% Allow extra inter-word stretch in paragraphs containing long unbreakable
% \texttt{} tokens (commit hashes, env-var names, parameter names) so they
% do not produce overfull hboxes.
\emergencystretch=3em
\sloppy

\theoremstyle{plain}
\newtheorem{theorem}{Theorem}
\newtheorem{proposition}[theorem]{Proposition}
\newtheorem{lemma}[theorem]{Lemma}
\theoremstyle{definition}
\newtheorem{definition}[theorem]{Definition}
\newtheorem{remark}[theorem]{Remark}

\newcommand{\bpp}{\mathrm{bpp}}
\newcommand{\kl}{D_{\mathrm{KL}}}
\newcommand{\Linears}{\mathcal{L}}
\newcommand{\Formats}{\mathcal{F}}
\newcommand{\Assignment}{\mathcal{A}}
\newcommand{\R}{\mathbb{R}}
\newcommand{\E}{\mathbb{E}}

\title{\bfseries PrismaQuant and PrismaSCOUT:\\
       Operational Rate--Distortion for\\
       Mixed-Precision Large Language Model Quantization}

\author{Robert Tand\\
        \textit{Independent Researcher}\\
        \texttt{robert.tand@icloud.com}}

\date{May 2026}

\begin{document}
\maketitle

\begin{abstract}
We describe \textsc{PrismaQuant}, a mixed-precision quantization pipeline for
large language models, and \textsc{PrismaSCOUT}, the search-and-select algorithm
at its core. Mixed-precision allocation is conventionally posed as a
constrained discrete optimization in which a per-Linear surrogate cost
$\sum_i u_i(f_i)$ is minimized over format assignments under a total-bit budget.
We argue that this formulation is structurally biased on modern transformers:
the true joint quantization error is non-additive across layers, and surrogate
ranks reliably overshoot when many flips are committed at once. We therefore
decouple \emph{candidate generation} from \emph{candidate selection}. Surrogate
costs propose discrete allocations cheaply via a multi-level hierarchy
($L_1$ probe, $L_2$ perturbed-activation fixed point, $L_3$ propagated end-KL);
every assignment that may ship is then gated by a real, end-to-end
Kullback--Leibler divergence measured on a held-out calibration split. Selection
is performed by a validated-frontier kneedle: the bisection samples
$(\bpp,\,\mathrm{KL})$ pairs adaptively, the dominance filter is widened by an
explicit KL noise floor, and the chosen knee is verified by a
leave-one-out stability diagnostic. A monotone coordinate-descent polish
guarantees the shipped assignment is no worse than the chosen frontier point.
On Qwen3.6\nobreakdash-27B, the pipeline produces a 5.31 bits-per-parameter
artifact with held-out KL $0.0151$, replacing a previously shipped 5.50 bpp
artifact at KL $0.0475$ from the same source weights and the same
export-time quantization tricks --- 11\% smaller and 68\% lower KL on the same
calibration. The artifact is published with DOI \texttt{10.57967/hf/8656}
\cite{Tand2026hf}. We also document, in a separate methodology section, several
numerical methods that were proposed and either implemented-but-deprecated or
rejected outright during the design process: Lagrangian
$\lambda$-bisection on the dual gradient, sandwich proximal recalibration,
block-DP over architectural cliques, sparse pairwise QUBO refinement, top-$K$
Hessian-eigenvalue covering, surrogate-only knee selection, and a
single-anchor probe-only knee predictor. We give the principled argument for
each, and the empirical or structural reason it did not become the shipping
path.
\end{abstract}

\section{Introduction}

Post-training quantization has converged on two dominant axes. Per-tensor
methods --- GPTQ \cite{gptq2022}, AWQ \cite{awq2024}, AutoRound
\cite{autoround2023}, HALO/HQQ-style rotation preprocessors --- improve the
per-Linear quantization error of a fixed format. Allocation methods --- HAWQ-V3
\cite{hawqv3}, CLADO \cite{clado2023}, and the more recent ParoQuant
\cite{paroquant}, IMPQ \cite{impq}, and AMQ \cite{amq} --- choose how many bits
to spend where. The two axes are largely orthogonal: an allocator can sit on top
of any per-tensor toolkit and decide which Linears get 16, 8, 4, 3, or 2 bits.

The allocation literature is unified by a recurring observation: per-layer
sensitivity surrogates (diagonal Hessian traces, Fisher MSEs, output-MSEs) are
biased estimators of joint quantization error. The standard remedy is to model
the bias --- pairwise ADMM in CLADO, sensitivity-matrix integer programming in
HAWQ, structured Pareto solvers in ParoQuant. We take a different tack.
\emph{Surrogates generate, real measurement selects.} An allocator does not need
a perfect cost model if every candidate it proposes can be cheaply re-scored
end-to-end on a held-out calibration split. Cross-layer interactions and
propagated effects then cease to be quantities one must \emph{model}; they are
quantities one \emph{observes}, automatically, every time a candidate is gated.

This paper reports the design and engineering of a system, named
\textsc{PrismaQuant}, that operationalizes that principle, and the search
algorithm at its core, \textsc{PrismaSCOUT} (\textbf{S}urrogate-\textbf{C}ascaded
\textbf{O}ptimization \textbf{U}nder \textbf{T}radeoff). Three contributions are
new:

\begin{enumerate}[leftmargin=*,itemsep=2pt]
  \item A \textbf{multi-level cost hierarchy} ($L_1, L_2, L_3$) that separates
    cheap proposal from increasingly faithful measurement. $L_2$ is a
    perturbed-activation fixed point that folds inter-layer coupling into the
    activation cache rather than into a closed-form interaction term. $L_3$ is
    a propagated end-KL measurement performed against a target-specific BF16
    paired baseline, lane-batched and tail-amortized so that one $L_3$ pass at
    Qwen-27B scale takes a tractable wall time.
  \item A \textbf{validated-frontier kneedle}: the elbow of the operational
    rate--distortion curve is found on \emph{measured} $(\bpp,\,\mathrm{KL})$
    pairs, not on a surrogate hull. We add a KL noise floor to the dominance
    filter, an adaptive bisection that refines around the current best knee, a
    leave-one-out stability diagnostic, and a strict monotone polish.
  \item A documented \textbf{methodology of rejected detours}. Several
    numerically attractive ideas (sandwich recalibration, $\lambda$-bisection,
    block-DP, sparse pairwise QUBO, top-$K$ Hessian covering, surrogate-only
    knee, probe-only prediction) were proposed during design, some implemented;
    we report the conditions under which each is principled and the empirical
    reasons each was not adopted as the shipping path.
\end{enumerate}

\noindent\textbf{Headline empirical result.} On Qwen3.6\nobreakdash-27B, with
identical source weights and identical export-time per-tensor tricks (HALO,
GPTQ damp sweep, scale sweep, block-output match), \textsc{PrismaSCOUT} produces
a 5.31~bpp artifact with held-out KL $0.0151$. The previously published
\textsc{PrismaQuant}~v1 artifact at the same source was 5.50~bpp with held-out
KL $0.0475$ from the same calibration. The new artifact is 11\% smaller and the
held-out KL is 68\% lower. Only the selection routine changed.

The paper is organized as follows. Section~\ref{sec:related} gives background
and related work. Section~\ref{sec:theory} sets up the theoretical framework:
the rate--distortion view, the additive surrogate failure mode, the cost
hierarchy, the kneedle. Section~\ref{sec:algorithm} states the algorithm.
Section~\ref{sec:rejected} is a deliberate methodology section on the numerical
detours we considered and rejected. Section~\ref{sec:engineering} reports the
engineering work that made the measurement loop tractable.
Section~\ref{sec:results} reports empirical results on Qwen3.6-4B and
Qwen3.6-27B, including an honest accounting of the tool-eval-bench delta.
Section~\ref{sec:discussion} discusses limitations.

\section{Background and Related Work}\label{sec:related}

\paragraph{Per-tensor quantization.} GPTQ \cite{gptq2022} is the workhorse
weight-only quantizer for transformers, posed as Hessian-weighted least squares
over reconstruction error. AWQ \cite{awq2024} reweights salient channels.
SmoothQuant \cite{smoothquant} pushes activation quantization tractable by
migrating activation magnitude into weights. AutoRound \cite{autoround2023}
sign-rounds via signSGD. HALO/HQQ-style rotations preprocess weight matrices to
reduce outlier counts before integer quantization. These methods take a single
matrix and a fixed format and produce a high-quality quantized version of that
matrix; they do not decide \emph{which} matrix gets which format.

\paragraph{Mixed-precision allocation.} HAWQ-V3 \cite{hawqv3} popularized
Hessian-eigenvalue-driven sensitivity allocation: layers with larger top
eigenvalues receive more bits. CLADO \cite{clado2023} extends to cross-layer
error coupling via ADMM on a structured pairwise cost. Recent work
(IMPQ \cite{impq}, ParoQuant \cite{paroquant}, AMQ \cite{amq}) frames the
problem in increasingly explicit Pareto terms, generally over surrogate cost
models. Surveys of the broader space have appeared as the field has matured
\cite{mpq_survey,iterquant2026}. The geometry-aware quantization line
\cite{geometry_aware} views per-layer quantization error through Riemannian
manifold structure of weight space.

\paragraph{Knee detection and operational rate--distortion.} The kneedle
algorithm of \cite{kneedle2011} provides a non-parametric estimator of the
elbow of a tradeoff curve. The closely related operational
rate--distortion (ORD) framing of \cite{ortega1998} establishes that any
discrete-choice rate--distortion problem admits a Lagrangian relaxation whose
solutions trace the lower convex hull of the achievable points. Knees on the
ORD frontier correspond to phase transitions where additional rate buys
disproportionate fidelity. Our work imports the ORD vocabulary --- and avoids
the convex hull, since for mixed-precision LLM quantization the discrete
frontier admits non-convex segments and a $\lambda$-sweep on the dual would
skip them (Section~\ref{subsec:lambda}).

\paragraph{Microscaling formats.} NVFP4 and MXFP4 \cite{microscaling} package
4-bit weights with FP8 group scales (NVFP4 group size 16; MXFP4 group size 32).
NVFP4 weight-only Linears are 4.5 bpp ($4 + 8/16 = 4.5$); MXFP4 is 4.25 bpp.
Modern Blackwell/Hopper GPUs offer first-class kernels for both. Mixed
precision in this paper means an assignment $\Assignment : \Linears \to
\Formats$ where each Linear independently selects from
$\{$BF16, MXFP8, NVFP4, MXFP4, NVINT3, $\dots\}$, subject to fused-sibling
constraints in attention ($q,k,v$) and MoE ($\mathrm{gate\_up},\mathrm{down}$).

\paragraph{Serving.} The artifacts described here are exported to vLLM
\cite{vllm} via the \texttt{compressed-tensors} schema, with format coverage
extending to NVFP4, MXFP8, BF16 passthrough, and (in MoE settings) MXFP4
experts. End-to-end inference uses Blackwell-native kernels.

\section{Theoretical Framework}\label{sec:theory}

\subsection{The shipping criterion}

Let $p_\theta$ denote the full-precision teacher model with parameters $\theta$
and $p_{\hat\theta}$ a quantized variant with weights $\hat\theta = Q_\Assignment(\theta)$
under format assignment $\Assignment$. Given a calibration distribution
$\mathcal{D}$ over input prefixes $x_{1:t}$, the operational quality of
$\Assignment$ is

\begin{equation}\label{eq:kl}
  \kl(\Assignment) \;=\; \E_{x \sim \mathcal{D}}\,
   \kl\!\left[\,p_\theta(\,\cdot \mid x)\,\big\|\,p_{\hat\theta}(\,\cdot \mid x)\,\right].
\end{equation}

KL is the natural objective for ``does the quantized model behave like the
teacher?'' because it captures the entire output distribution and not merely
the top-1 token. Two assignments with identical greedy decoding can differ
substantially in KL --- and that difference dominates sampling, instruction
following, and tool use, all of which depend on probability mass at non-modal
tokens. Throughout this paper KL is measured token-by-token, masked to
non-padding positions, averaged across calibration, and computed on a
held-out calibration split that is not seen by the surrogate cost stages.

\subsection{The additive-surrogate failure mode}

A natural objective for allocation is the additive sum of per-Linear costs

\begin{equation}\label{eq:additive}
  U_{\mathrm{add}}(\Assignment) \;=\; \sum_{i \in \Linears} u_i(\Assignment_i)
\end{equation}

where $u_i(f)$ is the local cost of using format $f$ at Linear $i$, holding
all other Linears at some reference assignment $\Assignment^{(\mathrm{ref})}$.
The diagonal Hessian formulation
$u_i(f) = \mathrm{Tr}(H_i)\cdot \mathrm{MSE}_i(f)$ used by HAWQ
\cite{hawqv3} is one example; output-MSE under fixed activation context is
another. \emph{For any choice of $u_i$ that holds $\Assignment_{-i}$ fixed,
$U_{\mathrm{add}}$ is a first-order Taylor expansion of $\kl(\Assignment)$
centered at $\Assignment^{(\mathrm{ref})}$.} A constrained minimization of
$U_{\mathrm{add}}$ subject to a bit budget $B$ via additive multi-choice
knapsack DP greedily flips many Linears at once and routinely steps outside
the trust region of the linear approximation. Empirically, the proposed
allocation overshoots true KL by a factor of $1.3$--$1.5$ on the
configurations we measured; a 0.10 unit predicted decrease can produce a 0.08
unit measured \emph{increase}.

This is the structural hazard the rest of the paper is built around.
Figure~\ref{fig:nonadditivity} sketches the geometry.

\input{figures/fig_nonadditivity.tex}

\subsection{Surrogates generate, real measurement selects}

The remedy adopted in \textsc{PrismaSCOUT} is to use $U_{\mathrm{add}}$ purely
as a \emph{candidate generator} and to select from generated candidates by
real, end-to-end measurement of $\kl(\Assignment)$ on a held-out calibration
split. The objective inequality
\[
  U_{\mathrm{add}}(\Assignment) \le U_{\mathrm{add}}(\Assignment') \quad
  \not\Longrightarrow \quad \kl(\Assignment) \le \kl(\Assignment')
\]
fails on average for many-flip moves. The substitute property
\[
  \kl(\Assignment) \le \kl(\Assignment')\quad\text{(measured on held-out)}
\]
is observable for any pair we are willing to materialize and forward through
the model. The cost of \emph{making} that observation is what determines
whether the surrogates-generate paradigm is practical at LLM scale, which is
where the cost hierarchy enters.

\subsection{The cost hierarchy}\label{subsec:hierarchy}

\textsc{PrismaSCOUT} runs three escalating cost models on the same problem.

\begin{description}[leftmargin=2em,style=nextline]
  \item[$L_1$ probe ($\sim$\,seconds per Linear).] A coarse per-matrix
    surrogate. We use a Fisher-weighted MSE
    $u_i^{(1)}(f) = \mathrm{Tr}(\hat H_i)\cdot \mathrm{MSE}_i(f)$ where
    $\hat H_i$ is an empirical Fisher accumulated by a multi-chunk probe pass
    \cite{prismaquant}. The output is a budget-feasible assignment computed
    by additive multi-choice knapsack DP; this is the role of HAWQ-style
    sensitivity, used here only as a starting point.
  \item[$L_2$ perturbed-X fixed point ($\sim$\,minutes).] A more careful
    per-Linear measurement that captures the activation distribution
    \emph{under the current assignment} rather than under BF16. We install
    pre-quantization activation hooks on the model under $\Assignment^{(k)}$,
    forward calibration tokens, and cache the resulting activations
    $X^{(k)}$. We then measure $u_i^{(2)}(f) = \mathrm{MSE}\!\left(W_i X^{(k)}, Q_f(W_i) X^{(k)}\right)$
    for each (Linear, format) pair. The DP solves under $u^{(2)}$, producing
    $\Assignment^{(k+1)}$. The fixed-point iteration $\Assignment \to X \to
    u^{(2)} \to \Assignment$ converges in $2$--$3$ iterations on every model
    we have measured. Inter-layer coupling enters through the activation
    cache: late-layer activations under a low-bit early-layer assignment
    differ from BF16 activations, and the per-Linear cost adapts.
  \item[$L_3$ propagated cost ($\sim$\,tens of minutes).] An end-KL
    measurement against a \emph{target-specific BF16 paired baseline}. For
    each candidate $(\mathrm{Linear}\;\ell,\,\mathrm{format}\;f)$, we run two
    forwards: a baseline lane in which Linear $\ell$ is BF16 and every other
    Linear is at its $L_2$ assignment, and a candidate lane in which Linear
    $\ell$ uses format $f$. The cost is the end-KL between the candidate
    lane logits and the baseline lane logits. Lane batching across all
    candidates within a depth group amortizes the forward; tail-only
    measurement avoids re-running the head/embedding for each lane. The
    output is a per-(Linear, format) end-KL table $u^{(3)}_i(f)$ that the
    final DP consumes.
\end{description}

\input{figures/fig_cost_hierarchy.tex}

The three levels are wired in cascade: $L_1$ proposes, $L_2$ refines under the
proposed activation distribution, $L_3$ measures the end-effect locally
around the converged $L_2$ assignment.
$L_3$ semantics differ from $L_1$ and $L_2$ in a load-bearing way: $L_3$
\emph{is} an end-KL measurement, just one whose paired baseline freezes the
context outside the target Linear at the converged $L_2$ assignment rather
than at BF16. This subtraction makes $L_3$ a local end-KL, defensible as a
unary cost for a final DP.

\subsection{The validated-frontier kneedle}\label{subsec:kneedle}

After the cost hierarchy generates a set of candidate assignments
$\{\Assignment^{(j)}\}$ at multiple budgets, each candidate is materialized,
forwarded against the held-out calibration split, and assigned a measured pair
$(\bpp(\Assignment^{(j)}),\,\kl(\Assignment^{(j)}))$. The Pareto frontier of
these points is the set of measured candidates not dominated by any other
measured candidate. The dominance filter is widened by an explicit KL noise
floor $\eta$:

\begin{definition}[$\eta$-dominance]
$\Assignment$ $\eta$-dominates $\Assignment'$ if
$\bpp(\Assignment) \le \bpp(\Assignment')$ and
$\kl(\Assignment) + \eta < \kl(\Assignment')$.
\end{definition}

This avoids spurious points on the frontier created by KL measurement noise
within $\eta$ of a competing point. The kneedle is then computed on the
$\eta$-non-dominated frontier as the point of maximum perpendicular distance
to the secant joining the frontier endpoints, after standard min-max
normalization of both axes. Concretely, with normalized
$(\tilde x_i, \tilde y_i)$ on the frontier in increasing-$\bpp$ order, and
$\tilde y$ replaced by $1 - \tilde y$ so the curve is increasing,

\begin{equation}\label{eq:kneedle}
  i^\star \;=\; \arg\max_i\;
  \frac{\big| (\tilde y_n - \tilde y_1)\tilde x_i - (\tilde x_n - \tilde x_1)\tilde y_i + \tilde x_n \tilde y_1 - \tilde y_n \tilde x_1 \big|}{\sqrt{(\tilde y_n - \tilde y_1)^2 + (\tilde x_n - \tilde x_1)^2}}.
\end{equation}

The kneedle is endpoint-fallback: if the maximum-perpendicular point is one of
the two endpoints (a near-linear frontier), the central point is used instead
and a warning is emitted. Adaptive bisection refines the frontier by
evaluating the midpoint of the largest gap adjacent to the current best knee
until either $|\bpp_R - \bpp_L| < \tau$ or the evaluation budget is reached.

\input{figures/fig_kneedle.tex}

\subsection{Leave-one-out stability}\label{subsec:loo}

A frontier with $5$--$7$ points is small enough that the chosen knee can be
brittle. We add a leave-one-out diagnostic: for each frontier point $i$,
remove it and re-run the kneedle (Figure~\ref{fig:loo_stability}). If the
maximum bpp shift across removals exceeds tolerance, or the maximum KL
shift exceeds $\eta$, the search emits an instability warning, and the
artifact metadata records both the chosen knee and a guarded fallback
(a geometric-median point on the frontier). We do not hard-fail on
instability; we record it. The diagnostic catches configurations in which
one or two frontier points dominate the elbow geometry --- usually a sign
that the bracket is too wide or the frontier is under-sampled near the
knee.

\input{figures/fig_loo_stability.tex}

\subsection{Coordinate descent as monotone polish}\label{subsec:coord}

\begin{proposition}[Monotone non-regression]
Let $\Assignment^{(0)}$ be the chosen knee assignment with measured KL
$\kappa_0 = \kl(\Assignment^{(0)})$. The coordinate-descent polish iterates
single-Linear flips $\Assignment^{(t+1)} = \Assignment^{(t)}\!\big[i \to
f^\star\big]$ and accepts only flips that strictly satisfy
$\kl(\Assignment^{(t+1)}) < \kl(\Assignment^{(t)})$ measured on the held-out
calibration split. The output assignment $\Assignment^{(\infty)}$ satisfies
$\kl(\Assignment^{(\infty)}) \le \kappa_0$.
\end{proposition}

\begin{proof}
Each accepted flip strictly decreases the held-out KL on the same
calibration split; rejected flips leave it unchanged. The sequence
$\{\kl(\Assignment^{(t)})\}$ is monotone non-increasing.
\end{proof}

The proposition is trivial; the operationally important point is that the
polish does not need a cost model. It does not need to predict the change in
KL; it measures it. The flip ranking comes from $L_3$ unary costs (cheap), the
acceptance criterion is real KL on held-out (slow but rare). On Qwen3.6-4B at
the 4.5~bpp anchor, coord descent committed 6 flips out of 101 attempts and
moved KL from $0.371$ (post-$L_2$) to $0.245$, a $34\%$ improvement. None of
those flips would have been picked by the additive surrogate; many were
rejected by the surrogate's predicted-cost ranking but accepted by real KL.

\input{figures/fig_coord_descent.tex}

\section{The Algorithm: \textsc{PrismaSCOUT}}\label{sec:algorithm}

\subsection{Pipeline}

For each anchor budget $B$:
\begin{enumerate}[leftmargin=*,itemsep=2pt]
  \item \textbf{Probe ($L_1$).} Fisher-weighted MSE per (Linear, format).
    Solve additive DP at $B$ to get $\Assignment^{(0)}$.
  \item \textbf{Perturbed-X fixed point ($L_2$).} Iterate $k = 0,\dots,K$:
    install activation hooks under $\Assignment^{(k)}$, cache calibration
    activations, measure $u^{(2)}_i(f)$ for all (Linear, format), solve
    weighted-Hamming-capped DP at $B$ for $\Assignment^{(k+1)}$. Stop on
    weighted-Hamming convergence ($K \le 3$ in practice).
  \item \textbf{Propagated cost ($L_3$).} Select a bounded neighborhood of
    Linears (uncertain choices, non-cheapest current picks, safety set);
    for each, measure $u^{(3)}_i(f)$ via paired BF16/candidate lanes,
    lane-batched, tail-only. Solve frozen DP over the neighborhood at $B$.
  \item \textbf{Real-KL gate.} Materialize the proposed assignment, forward
    against the held-out calibration split, record measured KL. If KL
    regresses beyond $\rho \cdot \kappa_{\mathrm{before}}$ (default $\rho =
    0$), roll back to the prior assignment.
\end{enumerate}

For the validated-frontier kneedle: repeat the above for a small set of
anchor budgets, validate the resulting candidates on the held-out split,
filter by $\eta$-dominance, and select via Equation~\eqref{eq:kneedle} with
adaptive bisection. Run leave-one-out (Section~\ref{subsec:loo}). Run
coordinate-descent polish (Section~\ref{subsec:coord}). Record audit trail.

\input{figures/fig_pipeline.tex}

\subsection{The DP and fused-sibling constraints}

The allocator solves an additive multi-choice knapsack DP whose
\emph{units} are not individual Linears but \emph{refinement units} that group
fused siblings. Attention $q,k,v$ projections share a per-tensor global scale
in vLLM's fused kernel \cite{vllm}; in MoE, $\mathrm{gate\_up}$ and
$\mathrm{down}$ packed Linears must share schemes within an expert column.
The DP enforces these constraints natively by enumerating only joint
sibling-coherent format choices per unit. Fused-sibling aggregation uses
sum (joint-cost) for activation-coupled siblings and max for
isolated-sensitivity siblings; we report the rationale and an empirical
comparison in Appendix~\ref{app:fused}.

\subsection{Calibration disjointness}\label{subsec:disjoint}

\textsc{PrismaSCOUT} maintains two disjoint calibration splits:

\begin{itemize}[leftmargin=*,itemsep=2pt]
  \item \textbf{Train split} (\texttt{l3\_calib\_split=train}, default).
    Used for $L_2$ activation caching and $L_3$ paired-baseline measurement.
  \item \textbf{Validation split} (\texttt{kl\_calib\_split=validation},
    default). Used for the real-KL gate, the validated-frontier kneedle, the
    coordinate-descent polish, and the artifact metadata. Never seen by any
    cost-generation stage.
\end{itemize}

The renaming and split discipline was added in commit \texttt{2e670b3}
(``Implement PR-Polish-1: Rename validation\_kl and secure surrogate
knee'') after an audit identified that the original
\texttt{validation\_kl} field was, in fact, an in-sample number on the same
calibration split that drove $L_2$/$L_3$ cost generation. Until that fix
landed, \texttt{validation\_kl} could be reasonably misread as a held-out
metric; it was not.

\input{figures/fig_calib_split.tex}

\section{Methods Considered and Rejected}\label{sec:rejected}

The shipping algorithm above is the result of a deliberate triangulation
between several principled alternatives. We describe each, the math behind
it, and the empirical or structural reason it was not adopted as the default.
The intent is not exhaustiveness for its own sake --- it is to make explicit
that simpler, mathematically attractive answers were considered first.

\subsection{Lagrangian \texorpdfstring{$\lambda$}{lambda}-bisection on the dual gradient}\label{subsec:lambda}

\paragraph{The proposal.} The constrained problem
$\min_\Assignment U(\Assignment)\;\text{s.t.}\;B(\Assignment) \le B$ admits the
Lagrangian relaxation $L(\Assignment, \lambda) = U(\Assignment) + \lambda(B(\Assignment) - B)$.
Sweeping $\lambda$ traces the lower convex envelope of the achievable
$(\bpp,\,U)$ frontier. For each $\lambda$, the per-Linear minimizer
decouples: $\Assignment_i^\star(\lambda) = \arg\min_f\,u_i(f) + \lambda\,
b_i(f)$ where $b_i(f)$ is the bit count for Linear $i$ at format $f$. This
costs CPU seconds. By Danskin's theorem, the achieved budget
$B^\star(\lambda)$ is the subgradient of the dual function. The kneedle is
the region of maximum $|dB^\star/d\lambda|$ on log-$\lambda$. A few-step
bisection on $\lambda$ produces the kneedle without sweeping the full grid.

\paragraph{The math is correct.} For a strictly convex frontier, the proposal
is optimal: $5$--$7$ $\lambda$ evaluations isolate the kneedle to within
$0.04$ bpp.

\paragraph{Why we did not adopt it as the default.} Two reasons.

First, the discrete frontier is not strictly convex. Sub-bit combinatorial
packing produces concave pockets where a slightly suboptimal assignment lies
\emph{above} the convex hull. No strictly positive $\lambda$ selects those
points as global minimizers of $L$; sweeping $\lambda$ leaps over them. For
the macroscopic kneedle this is harmless --- the macroscopic curve is
overwhelmingly convex --- but our production data on Qwen3.6-27B exhibits a
non-trivial gap between the convex-hull knee at 5.857~bpp / KL~0.056 and the
true measured knee at 5.31~bpp / KL~0.015 (Figure~\ref{fig:lambda_skip}).
The discrete archive search and the validated-frontier kneedle both find the
5.31 point; a pure $\lambda$-sweep on the surrogate does not.

Second, the surrogate $U_{\mathrm{add}}$ used inside the Lagrangian is the
same additive quantity whose joint behavior overshoots true KL by 30--50\%.
Even if the dual perfectly traced the surrogate hull, the surrogate hull is
not the true KL hull. We therefore retained a \emph{$\lambda$-archive search}
as an experimental candidate generator (commit \texttt{86f2c14}, ``Add
lambda archive knee search''; commit \texttt{444bc28}, ``Add one-pass Pareto
archive DP''), with a beam variant retaining alternate states per bit bin.
But the final selector is real-KL on the validated frontier; the
$\lambda$-archive feeds candidates into that selector, not the other way
around.

\input{figures/fig_lambda_sweep.tex}

\subsection{Sandwich proximal recalibration}

\paragraph{The proposal.} The non-additivity bug is a Taylor-trust-region
bug: $u^{(2)}_i$ is measured holding $\Assignment_{-i}$ at $\Assignment^{(\mathrm{ref})}$,
the DP then proposes a multi-flip $\Assignment^{(\mathrm{prop})}$ far outside
that radius. The remedy is a \emph{sandwich}: propose
$\Assignment^{(\mathrm{prop})}$ from the additive DP, re-measure
$u^{(2)}_i$ holding $\Assignment_{-i}$ at $\Assignment^{(\mathrm{prop})}$,
re-solve the DP. By Mirror-Descent convergence, the contraction factor
$\rho$ is the ratio of off-diagonal Hessian energy to block-diagonal
energy (empirically $\rho \in [0.3, 0.5]$). $2$--$3$ sandwich passes reach
within $5\%$ of the true KL minimum.

\paragraph{Why not adopted as default.} Two practical reasons.

First, each sandwich pass costs another full $L_2$ activation cache plus a
full DP solve. On Qwen3.6-27B, each pass is $\sim$22 minutes wall, so 3
passes is over an hour per anchor. A coordinate-descent polish of the same
anchor finds the same KL improvement at a fraction of the cost, because it
spends measurement on flips that are already small (single-Linear) and in a
region where the linear approximation does hold. Coord descent is also
trivially monotone on real KL by construction, while sandwich is
\emph{empirically} contractive.

Second, the operational re-centering target is the \emph{actual} chosen
assignment, not a sequence of $L_2$ DP outputs. The validated-frontier
kneedle and the coord-descent polish do this re-centering directly: the
gate is real KL on the assignment we are about to ship, with no Taylor
hand-waving in between. Sandwich would still be valuable as a refinement of
$L_2$ specifically; we list it under future work.

\subsection{Block-DP over architectural cliques}

\paragraph{The proposal.} The transformer Hessian is not dense; it is
block-diagonal with strong sequential off-diagonals confined to within a
single attention or MLP block. The structurally non-zero pairwise
interactions are essentially:

\[
  \mathrm{V}_{\mathrm{proj}} \to \mathrm{O}_{\mathrm{proj}}, \qquad
  \mathrm{Gate}_{\mathrm{proj}}/\mathrm{Up}_{\mathrm{proj}} \to \mathrm{Down}_{\mathrm{proj}}.
\]

A block-DP over these cliques (size $2$--$3$, disjoint after LayerNorm
resets, illustrated in Figure~\ref{fig:smrf_clique}) captures the dominant
interaction energy without an NP-hard QUBO, in
$K = 2 \cdot \mathrm{num\_layers}$ measurements per pass.

\input{figures/fig_smrf_clique.tex}

\paragraph{Why not adopted as default.} The shipping path already ships
fused-sibling aggregation in the additive DP, which bakes the dominant
intra-block coupling into joint sibling-coherent format choices. The
$\mathrm{Gate}/\mathrm{Up}\to\mathrm{Down}$ correlation is already
treated jointly inside the refinement unit. Block-DP would extend this to
$\mathrm{V}\to\mathrm{O}$, which is real, but not large enough on the
configurations we measured to justify a measurement pass beyond the existing
sibling aggregation. We retain it as a candidate for stronger interaction
regimes (sub-3-bit, MXFP4-only experts, reduced-rotation pre-processing).

\subsection{Sparse pairwise QUBO / SMRF}

\paragraph{The proposal.} For a sparse set of high-impact unit pairs $(i,j)$,
measure the residual
$\mathrm{res}_{ij}(a,b) = \mathrm{KL}_{\mathrm{paired}}(i=a,j=b) -
u_i(a) - u_j(b)$
and solve the resulting sparse QUBO in addition to the additive DP. This is
implemented in \texttt{interaction\_refine.py}.

\paragraph{Why default-off.} On the production Qwen3.6-27B run, the pairwise
stage covered 8 high-impact units at the cost of $\sim$250 paired forwards
($\sim$45 minutes) and rolled back when validation regressed. Coverage of
8 units out of $\sim$500 Linears is what one of our reviewers (Gemini~3.1)
called \emph{homeopathic}: too local to fix global non-additivity and
expensive enough to dominate the budget when activated. The audit consensus
hid the QUBO alias from the user-facing CLI, time-capped each pass, and
required a real-KL validation gate around any proposal it produces. The
code stays available for experimentation; the shipping path does not enable
it.

\subsection{Top-\texorpdfstring{$K$}{K} Hessian-eigenvalue covering}

\paragraph{The proposal.} Allocate measurement budget to the $K$ Linears
with largest top Hessian eigenvalue.

\paragraph{Why rejected.} The eigenvalue ranks individual matrices by local
sensitivity. It is structurally blind to the propagation graph: a Linear with
small top eigenvalue but a long downstream path through unnormalized
attention can dominate end-KL. The sandwich and validated-frontier results
both repeatedly select Linears that the eigenvalue ranking would have
missed.

\subsection{Surrogate-only knee selection}\label{subsec:surrogate}

\paragraph{The proposal.} Sweep target $\bpp$ on a single cached $L_3$ cost
table; let $U_{\mathrm{add}} = \sum_i u^{(3)}_i(\Assignment_i)$ stand in for
KL; pick the kneedle on the surrogate frontier.

\paragraph{Why rejected.} On the Qwen3.6-27B production data
(Figure~\ref{fig:surrogate_vs_validated}), the surrogate kneedle picks
5.857~bpp / KL~0.056. The validated kneedle on the same candidate set picks
5.31~bpp / KL~0.015. The surrogate is monotone within the additive trust
region; outside it, the ordering on $\bpp$ is no longer informative about
the ordering on KL. We retain the surrogate sweep as an
extremely cheap candidate generator (CPU seconds for hundreds of points)
and pair it with an obligatory \texttt{--knee-surrogate-validate} gate that
re-measures the picked points on real KL. After PR-Polish-1, the
surrogate-only path is hidden from the user-facing CLI; outputs are
relabeled \texttt{candidate\_surrogate\_*} so a reader cannot mistake the
surrogate kneedle for the shipping artifact.

\subsection{Probe-only knee prediction}

\paragraph{The proposal.} The kneedle is an artifact of intra-tensor
representation thresholds (the bpp at which a Linear's representation
capacity collapses). The $L_1$ probe captures those thresholds well. The
curvature of the probe sum and the curvature of the true frontier match up
to a multiplicative constant. Predict the knee bpp from $L_1$ alone.

\paragraph{Why used as a prior, not a selector.} On Qwen-4B-class scales the
expected error in the probe-predicted knee is empirically $\pm 0.12$--$0.18$
bpp \cite{prismaquant}. The density of mixed-precision states near the true
knee is often $\le 0.05$ bpp. A $\pm 0.15$ bpp error therefore guarantees we
miss the exact discrete kneedle. The probe is the optimal prior for
\emph{bracket} establishment (e.g., search in $[3.8, 4.1]$ bpp), but the
final discrete selection requires a real-KL evaluation of $3$--$5$
$\lambda$- or DP-generated candidates within that bracket.

\subsection{Pareto archive DP and beam variants}

\paragraph{The proposal.} A one-pass DP that retains, at each bit bin, the
Pareto-optimal set of $(\bpp,\,U)$ states. The archive's frontier becomes a
ready-made surrogate Pareto curve; the validated-frontier kneedle picks from
it. The beam variant
(\texttt{-{}-knee-archive-beam-per-bin}) retains $k$ alternate states per
bin. The refinement stage adds neighbors of the current best measured KL.

\paragraph{Status.} Implemented (commit \texttt{444bc28}; beam in commit
\texttt{53fe0db}). Used as a candidate generator in the production
Qwen3.6-27B archive run (\texttt{27b-knee-archive-beam4-refine-noseed}); the
exact 5.27~bpp/KL~0.013 supplied-best assignment was not reproduced by the
no-seed archive --- the no-seed selected $5.31$~bpp / KL~$0.0147$, about
9.9\% relative KL worse. We treat archive DP as a high-coverage candidate
generator that requires real-KL validation to ship, and not as a stand-alone
selector.

\input{figures/fig_method_decision_tree.tex}

\section{Engineering Considerations}\label{sec:engineering}

The theoretical framework requires real-KL measurement to dominate
selection. At LLM scale, this is only practical with attention to memory,
launch overhead, and determinism. We report the parts of the engineering
work that materially shape the algorithm.

\subsection{The 121~GB UMA constraint}

The development host is an NVIDIA GB10 (Blackwell) with 121~GB unified
memory shared between GPU and host. Quantization workloads at Qwen-27B scale
have a non-trivial peak memory schedule
(Figure~\ref{fig:memory_landscape}):
model weights $\sim$54~GB BF16, $L_2$ activation cache up to 24~GB across
iterations, frozen-weight cache $\sim$4~GB, lane-batched activations
$\sim$3$-$5~GB, CUDA graph pools $\sim$10$-$20~GB if naively stacked,
Triton kernel JIT transients $\sim$1$-$2~GB, and Python plus framework
baseline $\sim$3~GB. Single-component runs fit cleanly. All-on stacks reach
99--121~GB and OOM-kill, sometimes silently before the in-process budget
enforcer can intervene.

The empirical lesson, reflected throughout the codebase, is that the
shipping algorithm runs on the eager, default-pool path with the activation
cache and frozen-weight cache as the only large-state components, and that
all advanced kernels and graph optimizations are env-gated and disabled by
default. A run that takes longer but completes with a real-KL gate is
preferred over a run that benchmarks faster on a single pass and OOMs out at
coord descent.

\input{figures/fig_memory_landscape.tex}

\subsection{Lane-batched $L_3$ measurement}

The $L_3$ measurement of one (Linear, format) candidate forwards a calibration
batch through the model with one Linear in BF16 and one Linear at the
candidate format. Naively, each candidate is a separate forward. Lane
batching groups all candidates within a depth group into a single forward
with one extra batch dimension; per-lane Linear behavior is dispatched in
the swap hooks (the target's paired baseline lane stays BF16, the candidate
lane applies the candidate format, lanes belonging to other targets in the
same depth group apply that module's $L_2$ format). This reduced launch
overhead by 30$-$60$\times$ versus the sequential implementation, at the
cost of activation memory linear in the number of lanes; the
\texttt{-{}-l3-max-lanes-per-batch} flag bounds memory.

The tail-only mode caches the inputs to the tail of the network and replays
only the tail under each candidate, saving the head/embedding cost.
\texttt{PRISMAQUANT\_L3\_MIN\_HOST\_MEM\_GB} adds memory backpressure for
UMA hosts: the loop checks \texttt{/proc/meminfo} \texttt{MemAvailable}
between paired-override chunks and reduces the next chunk's lane count
before the host reaches OOM pressure.

\subsection{CUDA graphs and the shared-pool retreat}

CUDA graph capture eliminates per-launch overhead for the fixed-shape forward
paths used in $L_2$, $L_3$, KL evaluation, and coordinate-descent flip
evaluation. Naive capture allocates a separate graph memory pool per path, so
$4$--$5$ pools stack to $\sim$10$-$20~GB. We attempted a shared-pool
implementation (commit \texttt{3f320c7}) that captures all paths into a
single workspace with a tracked allocator. The shared-pool implementation
exposed a smoke-test NaN under one configuration (commit \texttt{d830a01},
``graphs: default shared pool OFF after smoke NaN risk''); a second pass
re-enabled shared-pool by default after a determinism flag and a
safety-override no-clone path were added (commits \texttt{2ba832f} and
\texttt{d9d67a9}). Through this, every path remains env-gated:
\texttt{PRISMAQUANT\_L3\_CUDA\_GRAPHS},
\texttt{PRISMAQUANT\_COORD\_LANE\_CUDA\_GRAPHS},
\texttt{PRISMAQUANT\_KL\_CUDA\_GRAPHS} accept \texttt{auto}/unset, \texttt{1}, or
\texttt{0}, and auto avoids graphing one-shot flip batches because capture
warmup dominates on small workloads.

The take-home is that CUDA graphs are real wall-time wins for stable-shape
loops (KL eval, $L_2$ activation cache, repeated lane forwards) and a real
hazard for exploratory loops (coord descent flip swaps). The shipping
default is auto: graph the stable paths, eager the rest.

\subsection{Determinism and reproducibility}

The audit trail recorded with each artifact includes: every measured
$(\bpp,\,\mathrm{KL})$ point on the validated frontier, the dominance filter
parameters ($\eta$ and the bisection tolerance), the chosen knee, the
leave-one-out shadow knees, the coord-descent flip log, the $L_2$ iteration
trace, the $L_3$ neighborhood selection rationale, and the final assignment
hash. The run is reproducible up to floating-point determinism on the
hardware that produced it (Blackwell GB10), which we have verified across
restarts on the same host. Artifacts are gated by an artifact registry
(\texttt{prismaquant.artifact\_registry}, commit \texttt{b3d8445}) that
refuses to publish without a passed validation entry.

\section{Empirical Results}\label{sec:results}

\subsection{Qwen3.6-4B at the 4.5~bpp anchor}

The 4.5~bpp anchor on Qwen3.6-4B is the canonical microbenchmark for the
non-additivity story. Figure~\ref{fig:4b_progression} reports KL at three
stages of the pipeline.

\input{figures/fig_4b_progression.tex}

The $L_3$ polish stage \emph{regresses} relative to $L_2$ ($0.371 \to 0.461$,
$+24\%$), consistent with the additive overshoot mechanism described in
Section~\ref{sec:theory}; the regression is detected by the real-KL gate
and the polish output is rolled back. The coordinate-descent stage
recovers and surpasses both: $0.461 \to 0.245$, with 6 flips committed out
of 101 attempts. The shipped 4B artifact at this anchor is the
coord-descent output. The same pattern reproduces at every Qwen-4B anchor we
have measured.

\subsection{Qwen3.6-27B production}\label{subsec:27b_results}

The Qwen3.6-27B PrismaSCOUT artifact (DOI \texttt{10.57967/hf/8656},
\cite{Tand2026hf}) is produced from the same source weights and
the same per-tensor toolkit
(HALO, GPTQ damp sweep, scale sweep, block-output match) as the previously
shipped \textsc{PrismaQuant}~v1 5.5~bpp artifact \cite{Tand2026hf_prev}. Only
the selection routine changed.

\begin{table}[h]
\centering
\caption{Qwen3.6-27B: previously shipped \textsc{PrismaQuant}~v1 vs.\ current
  \textsc{PrismaSCOUT} artifact, same source weights and same per-tensor
  toolkit. Held-out KL on disjoint validation split, 8 samples
  $\times$ 512 tokens.}
\label{tab:27b_headline}
\begin{tabular}{lrrr}
\toprule
Artifact & Size & bpp & Held-out KL \\
\midrule
\textsc{PrismaQuant}~v1 (5.50 bpp) & 22.67 GB & 5.50 & 0.0475 \\
\textsc{PrismaSCOUT} (5.31 bpp)    & 20.17 GB & 5.31 & 0.0151 \\
\midrule
Change                              & $-2.5$ GB ($-11\%$) & $-0.19$ & $-0.0324$ ($-68\%$) \\
\bottomrule
\end{tabular}
\end{table}

The artifact is 11\% smaller and the held-out KL is $\sim 3.2\times$ lower.
The validated-frontier picture (Figure~\ref{fig:validated_27b}) is the
algorithmic explanation: the surrogate-only kneedle on the same candidate
set picks 5.857~bpp / KL~0.056; the validated kneedle picks 5.31~bpp /
KL~0.015. The surrogate ranking on bpp does not preserve the ranking on
real KL inside the additive trust region's overshoot zone.

\input{figures/fig_validated_27b.tex}

\subsection{Downstream evaluation}

We report the downstream evaluation on the 27B PrismaSCOUT artifact in
Table~\ref{tab:27b_evals}. KL is a token-distribution statistic; it is not
a guarantee of task accuracy. We report results honestly, including a
specific tool-eval-bench regression we identified.

\begin{table}[h]
\centering
\caption{Qwen3.6-27B downstream evaluation under matched vLLM serve flags.
  Both artifacts served identically (\texttt{compressed-tensors},
  \texttt{kv-cache-dtype=fp8}, \texttt{enable-prefix-caching}). Validator and
  tool-eval-bench used \texttt{--no-think}.}
\label{tab:27b_evals}
\begin{tabular}{lrr}
\toprule
Metric & PrismaSCOUT (5.31) & PrismaQuant v1 (5.50) \\
\midrule
Held-out KL ($\downarrow$)               & 0.0151    & 0.0475    \\
Validator perplexity ($\downarrow$)       & 4.128     & 4.178     \\
Validator mean NLL ($\downarrow$)         & 1.4178    & 1.4297    \\
GSM8K strict ($\uparrow$)                  & 96.66\%   & --        \\
IFEval prompt strict ($\uparrow$)         & 85.40\%   & --        \\
MMLU 5-shot ($\uparrow$)                  & 86.23\%   & --        \\
tool-eval-bench hardmode ($\uparrow$)     & 88/100    & 85/100    \\
tool-eval safety-critical failures ($\downarrow$) & 4 & 3 \\
\bottomrule
\end{tabular}
\end{table}

The 5.31~bpp artifact wins on perplexity, mean NLL, and tool-eval-bench
hardmode score, but it has a stable regression on
tool-eval-bench scenario \texttt{TC-42} (extra parameter injection in tool
calls), and four versus three safety-critical failures. The shipped 5.5
artifact and the new 5.31 artifact share a vulnerability to scenario
\texttt{TC-60} (cross-turn sleeper injection); both follow an attacker BCC
inserted via a tool result. We do not claim KL improvement implies uniform
task improvement. Mean KL hides tail regressions and underweights
low-frequency schema and tool-call tokens. Future work includes p95/p99 KL
diagnostics and a schema-focused KL gate.

\section{Discussion and Limitations}\label{sec:discussion}

\paragraph{Non-additivity is structural, not pathological.} The 30--50\%
overshoot of the additive surrogate's predicted decrease versus measured KL
is not noise. It is the systematic bias of a first-order Taylor expansion
applied outside its trust region. The remedy is either a smaller step (which
$L_2$'s perturbed-X iteration approximates by re-centering the activation
cache) or a real-KL gate (which the validated-frontier kneedle and coord
descent enforce). The shipping path uses both.

\paragraph{The kneedle is a model-selection criterion.} The kneedle is not a
universal optimum. It selects the bpp where additional rate buys
disproportionately less fidelity --- i.e., where the marginal cost of further
size reduction is high. It is the right shipping criterion under the implicit
loss
$\mathcal{L}(\bpp,\,\mathrm{KL}) = \alpha\cdot\mathrm{KL} + \beta\cdot\bpp$
with $\alpha,\beta$ proportional to the perpendicular axis weights, and a
reasonable default in the absence of an explicit task-cost-of-bits curve.
Users with task-specific cost functions can substitute their own selector
on the validated frontier.

\paragraph{The held-out KL is a calibration-set statistic.} The downstream
evaluation in Section~\ref{subsec:27b_results} confirms that lower held-out
KL does not entail uniform task improvement. A schema-focused KL or a
top-token flip-rate diagnostic on chat-formatted prompts would tighten the
correspondence to tool-use accuracy. We report the KL we measured and flag
the gaps.

\paragraph{Hardware-bounded measurement.} The cost hierarchy is calibrated
to fit in a 121~GB UMA host; on larger hosts, $L_3$ candidate budgets and
coord-descent pass counts can be substantially larger, and the Pareto
sweep can target denser frontiers. The fundamental bound is that
real-KL measurement scales linearly with the number of validated assignments;
the validated-frontier kneedle keeps that count to $5$--$15$ per artifact.

\paragraph{Limitations.}
(i) The surrogate hierarchy is calibrated for transformer architectures;
its $L_2$ perturbed-X step assumes that activation caching is feasible, which
breaks for sparse-routing MoE without a careful expert-aware cache.
(ii) The sandwich, block-DP, and full QUBO are deliberately not adopted but
remain principled candidates for stronger interaction regimes.
(iii) Calibration is wikitext-derived and chat-formatted via a fixed system
prompt; domain-shifted deployments may benefit from in-domain calibration.
(iv) No claim of optimality is made beyond ``provably no worse than the
chosen frontier point under monotone real-KL flips.'' The frontier itself is
sampled, not exhaustive.

\section{Conclusion}

We described \textsc{PrismaQuant} and its core search algorithm
\textsc{PrismaSCOUT}: a mixed-precision LLM quantization pipeline whose
defining choice is to use surrogate cost models for candidate generation
only and to make every shipping decision on a real, end-to-end
Kullback--Leibler divergence measurement against a held-out calibration
split. The key engineering challenge was making that measurement loop
tractable on a 121~GB UMA host through a multi-level cost hierarchy with
lane-batched, tail-amortized $L_3$, validated-frontier kneedle selection
with a noise-floor dominance filter and leave-one-out stability, and a
monotone coord-descent polish that is provably non-regressive on real KL.

The shipping artifact on Qwen3.6-27B is 11\% smaller and 68\% lower in
held-out KL than the previously shipped \textsc{PrismaQuant}~v1 artifact at
the same source. We also documented several principled numerical methods
(Lagrangian $\lambda$-bisection, sandwich proximal recalibration, block-DP
over architectural cliques, sparse pairwise QUBO, top-$K$ Hessian covering,
surrogate-only knee, probe-only knee predictor) that were considered during
design and either deprecated, made experimental-only, or rejected outright.
Each is principled in its regime; none is the right shipping default for
mixed-precision LLM quantization under our measurement constraints.

\paragraph{Reproduction.} The Qwen3.6-27B PrismaSCOUT artifact is available at
\url{https://huggingface.co/rdtand/Qwen3.6-27B-PrismaSCOUT-Blackwell-NVFP4-BF16-vllm}
(DOI \texttt{10.57967/hf/8656}, \cite{Tand2026hf}). The pipeline is open
source at \url{https://github.com/RobTand/prismaquant}; commits referenced
throughout (\texttt{8e503a5}, \texttt{2e670b3}, \texttt{444bc28}, \dots) are
visible there.

\section*{Acknowledgments}

The author thanks two LLM collaborators (\textsc{Claude} and \textsc{Codex})
for sustained design and implementation help, and \textsc{Gemini~3.1~Pro}
for the mathematical formalism (Lagrangian dual analysis, sandwich
proximal-descent argument, block-DP clique structure, probe error bounds)
that produced the rejected-methods catalog of Section~\ref{sec:rejected}.
Where their contributions are load-bearing for the paper's narrative, the
relevant deliberation transcripts are preserved in
\texttt{scratch/powwow/} of the source repository for future audit.

\renewcommand{\refname}{References}

\begin{thebibliography}{99}

\bibitem{Tand2026hf} R.~Tand. \emph{Qwen3.6-27B-PrismaSCOUT-Blackwell-NVFP4-BF16-vllm.}
  Hugging Face model artifact, 2026. DOI: \texttt{10.57967/hf/8656}.
  \url{https://huggingface.co/rdtand/Qwen3.6-27B-PrismaSCOUT-Blackwell-NVFP4-BF16-vllm}.

\bibitem{Tand2026hf_prev} R.~Tand. \emph{Qwen3.6-27B-PrismaQuant-5.5bit-vllm.}
  Hugging Face model artifact, 2026.
  \url{https://huggingface.co/rdtand/Qwen3.6-27B-PrismaQuant-5.5bit-vllm}.

\bibitem{prismaquant} R.~Tand. \emph{PrismaQuant: per-Linear sensitivity-driven
  mixed-precision quantization for LLMs.}
  \url{https://github.com/RobTand/prismaquant}, 2026.

\bibitem{gptq2022} E.~Frantar, S.~Ashkboos, T.~Hoefler, and D.~Alistarh. GPTQ:
  Accurate post-training quantization for generative pre-trained transformers.
  \emph{ICLR}, 2023. \url{https://arxiv.org/abs/2210.17323}.

\bibitem{awq2024} J.~Lin, J.~Tang, H.~Tang, S.~Yang, X.~Dang, C.~Gan, and S.~Han.
  AWQ: Activation-aware weight quantization for LLM compression and
  acceleration. \emph{MLSys}, 2024. \url{https://arxiv.org/abs/2306.00978}.

\bibitem{autoround2023} W.~Cheng, Y.~Lv, W.~Zhang, and H.~Shen. Optimize weight
  rounding via signed gradient descent for the quantization of LLMs.
  \emph{arXiv:2309.05516}, 2023.

\bibitem{smoothquant} G.~Xiao, J.~Lin, M.~Seznec, H.~Wu, J.~Demouth, and S.~Han.
  SmoothQuant: Accurate and efficient post-training quantization for large
  language models. \emph{ICML}, 2023.

\bibitem{hawqv3} Z.~Yao, Z.~Dong, Z.~Zheng, A.~Gholami, J.~Yu, E.~Tan, L.~Wang,
  Q.~Huang, Y.~Wang, M.~W.~Mahoney, and K.~Keutzer. HAWQ-V3: Dyadic neural
  network quantization. \emph{ICML}, 2021. arXiv:2011.10680.

\bibitem{clado2023} D.~Lee, M.~Kim, and S.~Han. CLADO: Cross-layer ADMM
  optimization for joint mixed-precision quantization.
  \emph{arXiv:2307.05657}, 2023.

\bibitem{paroquant} (ParoQuant authors). Pareto-optimal mixed-precision
  quantization. \emph{arXiv:2511.10645}, 2025.

\bibitem{impq} (IMPQ authors). Integer-mixed-precision quantization for LLMs.
  \emph{arXiv:2509.15455}, 2025.

\bibitem{amq} (AMQ authors). Adaptive mixed-precision quantization.
  \emph{arXiv:2509.12019}, 2025.

\bibitem{mpq_survey} (Survey authors). Mixed-precision quantization for large
  language models: a survey. \emph{arXiv:2510.16805}, 2025.

\bibitem{iterquant2026} (IterQuant authors). Iterative mixed-precision
  quantization for LLMs. \emph{DATE}, 2026.

\bibitem{geometry_aware} (Authors). Geometry-aware quantization of large
  language models. \emph{arXiv:2507.18553}, 2025.

\bibitem{kneedle2011} V.~Satopaa, J.~Albrecht, D.~Irwin, and B.~Raghavan. Finding
  a ``kneedle'' in a haystack: Detecting knee points in system behavior.
  \emph{ICDCSW}, 2011.

\bibitem{ortega1998} A.~Ortega and K.~Ramchandran. Rate-distortion methods for
  image and video compression. \emph{IEEE Signal Processing Magazine},
  15(6):23--50, 1998.

\bibitem{microscaling} (OCP). Microscaling formats: MXFP4, MXFP6, MXFP8 and
  NVFP4. \emph{Open Compute Project Microscaling Spec.}, 2023--2024.

\bibitem{vllm} W.~Kwon, Z.~Li, S.~Zhuang, Y.~Sheng, L.~Zheng, C.~H.~Yu,
  J.~Gonzalez, H.~Zhang, and I.~Stoica. Efficient memory management for large
  language model serving with PagedAttention. \emph{SOSP}, 2023.

\bibitem{turboboa} (Authors). TurboBOA: Throughput-optimal mixed-precision
  inference. \emph{arXiv:2602.04929}, 2026.

\bibitem{shapley_potential_games} D.~Monderer and L.~Shapley. Potential games.
  \emph{Games and Economic Behavior}, 14:124--143, 1996.

\end{thebibliography}

\appendix

\section{Fused-sibling aggregation}\label{app:fused}

Fused siblings in attention ($q,k,v$) and MoE ($\mathrm{gate\_up}$,
$\mathrm{down}$) must share one format choice in vLLM's fused kernel. The
allocator's refinement unit groups them and enumerates only joint
sibling-coherent format combinations. The aggregation rule for the unit's
unary cost has two natural candidates: $\sum$ over siblings (sum-aggregation)
and $\max$ over siblings (max-aggregation). Sum is correct when the siblings'
output errors propagate independently; max is correct when one sibling
dominates the unit's downstream activation noise. Empirically on
Qwen3.6-27B, sum-aggregation matches the measured $L_3$ end-KL more closely
on both attention and MoE units (commit \texttt{d370c3e}, ``allocator: revert
to sum-aggregation for fused siblings (right thing on principle)''). We use
sum.

\section{Allocator runtime parameters}

\begin{itemize}[leftmargin=*]
  \item $L_2$: \texttt{bit\_precision}~$=0.001$, weighted-Hamming convergence
    threshold $0.05$, max iterations $5$ (typically converges in $2$--$3$).
  \item $L_3$:
    \texttt{l3-uncertainty-rel-tol}~$=0.10$,
    \texttt{l3-min-fraction}~$=0.05$,
    \texttt{l3-max-fraction}~$=0.30$,
    \texttt{l3-safety-fraction}~$=0.02$,
    \texttt{l3-max-lanes-per-batch}~$=64$ (UMA-cap aware),
    \texttt{l3-tail-only}~$=$~on,
    \texttt{l3-regression-tolerance}~$=0.0$.
  \item Validated-frontier kneedle: KL noise floor
    $\eta$ default $5\times 10^{-4}$, bisection tolerance $0.05$ bpp,
    leave-one-out at $\ge 4$ frontier points.
  \item Coord descent: \texttt{max-passes}~$=1$,
    \texttt{early-stop-streak}~$=50$,
    \texttt{max-lanes-per-batch}~$=64$.
\end{itemize}

\end{document}
